%% LyX 2.1.4 created this file.  For more info, see http://www.lyx.org/.
%% Do not edit unless you really know what you are doing.
\documentclass[a4paper,oneside,brazil,11pt,a4paper,openright,titlepage,usenames,dvipsnames]{book}
\usepackage[T1]{fontenc}
\usepackage[latin9]{inputenc}
\setcounter{secnumdepth}{3}
\setcounter{tocdepth}{3}
\usepackage{array}
\usepackage{verbatim}
\usepackage{calc}
\usepackage{textcomp}
\usepackage{amssymb}
\usepackage{graphicx}

\makeatletter

%%%%%%%%%%%%%%%%%%%%%%%%%%%%%% LyX specific LaTeX commands.
\pdfpageheight\paperheight
\pdfpagewidth\paperwidth

%% Because html converters don't know tabularnewline
\providecommand{\tabularnewline}{\\}

%%%%%%%%%%%%%%%%%%%%%%%%%%%%%% User specified LaTeX commands.
% Classe alternativa, apropriada para impressão frente-verso. Inclui páginas em branco
% de forma que capítulos sempre tenham início na página à direita:
% \documentclass[11pt,a4paper,openright,titlepage]{book}

% Pacotes
\usepackage[T1]{fontenc}
\usepackage[brazilian]{babel}
\usepackage{epsfig}
\usepackage{subfigure}
\usepackage{amsfonts}
\usepackage{amsmath}
\usepackage[thmmarks,amsmath]{ntheorem}%\usepackage{amsthm}
\usepackage{boxedminipage}
\usepackage{geometry}
\usepackage{theorem}
\usepackage{fancybox}
\usepackage{fancyhdr}
\usepackage{ifthen}
\usepackage{url}
\usepackage{afterpage}
\usepackage{color}
\usepackage{colortbl}
\usepackage{rotating}
\usepackage{makeidx}
\usepackage{indentfirst}
% Pacotes para adição de figuras do inkscape
\usepackage{graphicx}
\usepackage{import}

% Escolher um dos seguintes formatos:
\usepackage{ft2unb} % segue padrão de fontes do Latex

\makeindex

\makeatother

\usepackage{babel}
\begin{document}
\setcounter{secnumdepth}{3}
\setcounter{tocdepth}{2}
\pagestyle{empty}

\grau{Engenheiro de Controle e Automação}

\tipodemonografia{TRABALHO DE GRADUAÇÃO}

\begin{comment}
Título
\end{comment}


\titulolinhai{TÍTULO DO TRABALHO DIVIDIDO EM}

\titulolinhaii{MAIS DE UMA LINHA PARA TÍTULOS}

\titulolinhaiii{REALMENTE LONGOS COMO ESTE}

\titulolinhaiv{}

\begin{comment}
Autores. Basta retirar o texto totalmente caso não haja um determinado
autor.
\end{comment}


\autori{João da Silva}

\autorii{}

\autoriii{}

\begin{comment}
Membros da banca. Basta retirar o texto totalmente caso não haja um
determinado membro da banca.
\end{comment}


\membrodabancai{Prof. Fulano de Sousa, ENE/UnB}

\membrodabancaifuncao{Orientador}

\membrodabancaii{Prof. Ciclano dos Santos, ENE/UnB}

\membrodabancaiifuncao{Examinador interno}

\membrodabancaiii{Prof. Beltrano da Silva, ENE/UnB}

\membrodabancaiiifuncao{Examinador interno}

\membrodabancaiv{}

\membrodabancaivfuncao{}

\membrodabancav{}

\membrodabancavfuncao{}

\begin{comment}
Data de defesa: mês e ano
\end{comment}


\mes{dezembro}
\ano{2015}

\begin{comment}
Comandos para criar a capa e a página de assinaturas
\end{comment}


\capaprincipal
\capaassinaturas

\begin{comment}
Ficha Catalográfica
\end{comment}


\noindent \textbf{FICHA CATALOGRÁFICA}

\noindent %
\fbox{\begin{minipage}[t]{1\columnwidth}%
ARTHUR, DE MATOS BEGGS

RISC-V SiMPLE,

\medskip{}


{[}Distrito Federal{]} 2018.

\medskip{}


x, 101p., 297 mm (FT/UnB, Engenheiro, Controle e Automação, 2018).
Trabalho de Graduação \textendash{} Universidade de Brasília.Faculdade
de Tecnologia.

\medskip{}


1. RISC-V\hfill{}2. ???\hfill{}

\medskip{}


I. Mecatrônica/FT/UnB\hfill{}II. Título (Série)\hfill{}

%
\end{minipage}}

\noindent \medskip{}


\noindent \textbf{REFERÊNCIA BIBLIOGRÁFICA}

BEGGS, ARTHUR DE MATOS, (2018). RISC-V SiMPLE. Trabalho de Graduação
em Engenharia de Controle e Automação, Publicação FT.TG-$n^{\circ}022$,
Faculdade de Tecnologia, Universidade de Brasília, Brasília, DF, 101p.

\noindent \bigskip{}


\noindent \textbf{CESSÃO DE DIREITOS}

\noindent AUTOR: Arthur de Matos Beggs

TÍTULO DO TRABALHO DE GRADUAÇÃO: RISC-V SiMPLE.

\noindent \medskip{}


\noindent GRAU: Engenheiro\hfill{}ANO: 2018\hfill{}

\noindent \medskip{}


É concedida à Universidade de Brasília permissão para reproduzir cópias
deste Trabalho de Graduação e para emprestar ou vender tais cópias
somente para propósitos acadêmicos e científicos. O autor reserva
outros direitos de publicação e nenhuma parte desse Trabalho de Graduação
pode ser reproduzida sem autorização por escrito do autor.

\noindent \bigskip{}


\noindent \rule[0.5ex]{1\columnwidth}{1pt}

\noindent Arthur de Matos Beggs

\noindent SHCGN 703 Bl G Nº 120, Asa Norte

\noindent 70730-707 Brasília \textendash{} DF \textendash{} Brasil.


\begin{comment}
Dedicatória
\end{comment}


\frontmatter

\begin{comment}
Texto de dedicatória do primeiro autor.
\end{comment}


\dedicatoriaautori{um beijo pra minha mãe, pro meu pai, e pra você}

\begin{comment}
Texto de dedicatória do segundo autor. Caso não tenha um segundo autor,
este texto não será mostrado
\end{comment}


\dedicatoriaautorii{Dedicatória do autor 2}

\begin{comment}
Texto de dedicatória do terceiro autor. Caso não tenha um terceiro
autor, este texto não será mostrado
\end{comment}


\dedicatoriaautoriii{Dedicatória do autor 3}

\begin{comment}
Comando para criar a página de dedicatória
\end{comment}


\dedicatoria

\begin{comment}
Agradecimentos
\end{comment}


\begin{comment}
Texto de agradecimentos do primeiro autor.
\end{comment}


\agradecimentosautori{Agradecimentos!}

\begin{comment}
Texto de agradecimentos do segundo autor. Caso não tenha um segundo
autor, este texto não será mostrado.
\end{comment}


\agradecimentosautorii{A inclusão desta seção de agradecimentos é
opcional e fica à critério do(s) autor(es), que caso deseje(em) inclui-la
deverá(ão) utilizar este espaço, seguindo esta formatação.}

\begin{comment}
Texto de agradecimentos do terceiro autor. Caso não tenha um terceiro
autor, este texto não será mostrado.
\end{comment}


\agradecimentosautoriii{A inclusão desta seção de agradecimentos
é opcional e fica à critério do(s) autor(es), que caso deseje(em)
inclui-la deverá(ão) utilizar este espaço, seguindo esta formatação.}

\begin{comment}
Comando para criar a página de agradecimentos
\end{comment}


\agradecimentos

\resumo{resumo}{Resumo!

\medskip{}


Palavras Chave: bla, ble, bli

}\vspace*{2cm}


\resumo{Abstract}{Abstract, in English ofc!

\medskip{}


Keywords: bla, ble, bli

}

\begin{comment}
Listas de conteúdo, figuras e tabelas
\end{comment}


\sumario
\listadefiguras
\listadetabelas

\begin{comment}
Lista de Símbolos
\end{comment}


%TCIDATA{LaTeXparent=0,0,these.tex}


%\chapter*{\setfontarial\mdseries LISTA DE SÍMBOLOS} % se usar ft1unb.sty, descomente esta linha



\chapter*{LISTA DE SÍMBOLOS}

% se usar ft2unb.sty, descomente esta linha



\subsection*{Símbolos Latinos}

\begin{tabular}{p{0.1\textwidth}p{0.63\textwidth}>{\PreserveBacklash\raggedleft}p{0.15\textwidth}}
$v$  & Velocidade linear  & {[}m/s{]}\tabularnewline
\end{tabular}


\subsection*{Símbolos Gregos}

\begin{tabular}{p{0.1\textwidth}p{0.63\textwidth}>{\PreserveBacklash\raggedleft}p{0.15\textwidth}}
$\omega$ & Velocidade angular & {[}rad/s{]}\tabularnewline
\end{tabular}


\subsection*{Grupos Adimensionais}

\begin{tabular}{p{0.1\textwidth}p{0.8\textwidth}}
i, k & Contador\tabularnewline
\end{tabular}


\subsection*{Subscritos}

\begin{tabular}{p{0.1\textwidth}p{0.8\textwidth}}
$ref$  & referência \tabularnewline
$fer$  & ferramenta \tabularnewline
$sis$  & sistema \tabularnewline
$des$  & desejado\tabularnewline
\end{tabular}


\subsection*{Sobrescritos}

\begin{tabular}{p{0.1\textwidth}p{0.8\textwidth}}
$\cdot$  & Variação temporal \tabularnewline
$-$  & Valor médio \tabularnewline
\end{tabular}


\subsection*{Siglas}

\begin{tabular}{p{0.1\textwidth}p{0.8\textwidth}}
PCI  & \textit{Peripheral Component Interconnect}\tabularnewline
CPU & Unidade Central de Processamento - \textit{Central Processing Unit} \tabularnewline
AO & Saída Analógica - \textit{Analog Out}\tabularnewline
DO & Saída Digital - \textit{Digital Out}\tabularnewline
CS & Seletor de \textit{Chip - Chip Select}\tabularnewline
SC & Sem Conexão\tabularnewline
P.I. & Placa de Interface\tabularnewline
ICW & \textit{Initialization Command Words}\tabularnewline
OCW & \textit{Operational Control Word}\tabularnewline
\end{tabular}


\begin{comment}
Corpo Principal
\end{comment}


\mainmatter
\setcounter{page}{1}
\pagenumbering{arabic}
\pagestyle{plain}

\begin{comment}
Introdução
\end{comment}

\chapter{Introdução}

\label{CapIntro}

% Resumo opcional. Comentar se não usar.
\resumodocapitulo{Resumo opcional}


\section{Contextualização}

Contextualizar.

Conforme \cite{article:dummy}, vide a Tabela \ref{tab:Descrever-tabela}.
Assim sendo, observe a Figura \ref{fig:Descrever-figura.}.
\begin{table}[h]
\begin{centering}
\begin{tabular}{|c|c|c|c|c|}
\hline
 &  &  &  & \tabularnewline
\hline
\hline
 &  &  &  & \tabularnewline
\hline
 &  &  &  & \tabularnewline
\hline
 &  &  &  & \tabularnewline
\hline
 &  &  &  & \tabularnewline
\hline
\end{tabular}
\par\end{centering}

\label{tab:Descrever-tabela}Descrever tabela.


\end{table}


\begin{figure}[h]
\begin{centering}
\includegraphics[width=0.4\columnwidth]{figs/capa_fundo}
\par\end{centering}

\label{fig:Descrever-figura.}Descrever figura.


\end{figure}



\section{Definição do problema}

Definir problema.


\section{Objetivos do projeto}

Objetivos.


\section{Resultados obtidos}

Resultados.


\section{Apresentação do manuscrito}

Apresentar.


\begin{comment}
Fundamentos
\end{comment}

\chapter{Fundamentos\label{chap:FundamentacaoMatematica}}

% Resumo opcional. Comentar se n�o usar.
\resumodocapitulo{Resumo opcional.}


\section{Introdu��o}

Introduzir.


\begin{comment}
Conclusões
\end{comment}

\chapter{Conclusões}

\label{CapConclusoes}

Concluir


\section{Perspectivas Futuras}

Perspectivas futuras


\begin{comment}
Bibliografia
\end{comment}


\renewcommand{\bibname}{REFERÊNCIAS BIBLIOGRÁFICAS}
\addcontentsline{toc}{chapter}{REFERÊNCIAS BIBLIOGRÁFICAS}

\bibliographystyle{abnt-num}
\bibliography{bibliography}


\begin{comment}
Anexos
\end{comment}


\anexos
\makeatletter
% não retirar estes comandos
\renewcommand{\@makechapterhead}[1]{%
  {\parindent \z@ \raggedleft \setfontarial\bfseries
\LARGE \thechapter. \space\space
\uppercase{#1}\par
\vskip 40\p@
}
}
\makeatother

\begin{comment}
Anexo I: Descrição do CD
\end{comment}



\chapter{Descrição do conteúdo do CD}

\label{AnCD} 

Descrever CD.


\refstepcounter{noAnexo}

\begin{comment}
Anexo II: Programas Utilizados
\end{comment}



\chapter{Programas utilizados}

Quais programas foram utilizados?


\refstepcounter{noAnexo}

\begin{comment}
Acrescente mais anexos conforme julgar necessário.
\end{comment}

\end{document}
